\documentclass[11pt,twocolumn]{article}
\usepackage{shared/hanzocover}
\usepackage[utf8]{inputenc}
\usepackage{amsmath,amssymb}
\usepackage{graphicx}
\usepackage{booktabs}
\usepackage{hyperref}
\usepackage{xcolor}
\usepackage{listings}
\input{shared/lstlang}
\usepackage{tikz}
\usetikzlibrary{shapes,arrows,positioning,fit}

\definecolor{codegreen}{rgb}{0,0.6,0}
\definecolor{codegray}{rgb}{0.5,0.5,0.5}
\definecolor{codepurple}{rgb}{0.58,0,0.82}
\definecolor{backcolour}{rgb}{0.95,0.95,0.92}
\lstdefinestyle{mystyle}{backgroundcolor=\color{backcolour},commentstyle=\color{codegreen},keywordstyle=\color{codepurple},numberstyle=\tiny\color{codegray},stringstyle=\color{codegreen},basicstyle=\ttfamily\footnotesize,breaklines=true,captionpos=b,keepspaces=true,numbers=left,numbersep=5pt,showstringspaces=false,tabsize=2}
\lstset{style=mystyle}

\title{A Post-Quantum Object Store: Per-Tenant SQLite and a Dropbox-Class Filesystem over One Encrypted Block Substrate}
\author{Zach Kelling\thanks{zach@hanzo.ai} \\ \textit{Hanzo Industries} \\ \texttt{research@hanzo.ai}}
\date{June 2026}

\begin{document}
\hanzocoverpage

\begin{abstract}
We describe a single object store that presents two access surfaces over one
content-addressed, post-quantum--encrypted block substrate backed by
S3-compatible storage. The first surface is a native, in-process SQLite virtual
file system (\texttt{hanzovfs}) that gives every tenant a private, encrypted
database without FUSE or cgo. The second is a POSIX FUSE/WASI mount that exposes
the same substrate as an ordinary read/write directory---Dropbox-class file sync
with local-first caching. Both surfaces share one key hierarchy: per-object data
keys sealed with a hybrid ML-KEM-768 + X25519 KEM (\texttt{luxfi/age}) and
bulk encryption under ChaCha20-Poly1305. We show that the database surface
sustains 720K inserts/s and 59.6K point-reads/s---$2.6\times$ writes and
$5\times$ reads over a FUSE-mounted equivalent---while remaining end-to-end
post-quantum. This realizes HIP-0302 (encrypted SQLite durability) and HIP-0107
(streaming replication over VFS) as one deployable substrate.
\end{abstract}

\section{Introduction}
Multi-tenant platforms accrete storage engines: Postgres for rows, Redis for
cache, an object store for blobs, a separate sync service for files. Each adds a
network service to operate, a credential to rotate, and a blast radius. We argue
for the opposite: \emph{one} encrypted block substrate, addressed two ways.

A tenant's structured state lives in a private SQLite database; a tenant's files
live in a private directory; both are the same encrypted blocks in the same
bucket, distinguished only by the access surface mounted over them. There is no
second storage service, no shared relational database, and no plaintext at rest.
Crucially, the design is \emph{post-quantum by construction}: an adversary who
records every ciphertext object today cannot recover any tenant's data under a
future quantum computer.

\section{The Block Substrate}
The substrate stores fixed-size, individually-sealed blocks behind the
\texttt{Backend} interface (\texttt{Get/Put/Delete/List}), with concrete
backends for \texttt{file://} (local dev) and \texttt{s3://} (production,
served by \texttt{hanzoai/s3}). A local NVMe cache fronts the backend so hot
working sets never touch the network; cold blocks are fetched and decrypted on
demand. Objects are content-addressed: a block's key is derived from its sealed
bytes, which yields intra-tenant deduplication and cheap, immutable snapshots.
Cross-tenant deduplication is deliberately \emph{not} performed---it would leak
equality of plaintext across tenants---so each tenant's blocks are sealed under
that tenant's own key.

\section{Post-Quantum Key Hierarchy}
Symmetric encryption is already quantum-resistant: Grover's algorithm reduces a
256-bit key to a 128-bit security level, which remains infeasible. The quantum
threat is to \emph{asymmetric} primitives---Shor's algorithm breaks RSA, ECDH,
and EdDSA. We therefore place post-quantum cryptography exactly where it matters:
key encapsulation and signatures.

\textbf{Bulk.} Each block (or each 4\,KiB database page) is sealed with
ChaCha20-Poly1305 under a per-object \emph{data encryption key} (DEK), giving
confidentiality and integrity (AEAD) for the payload.

\textbf{Key wrap.} The DEK is sealed with \texttt{luxfi/age} using an HPKE
hybrid recipient, \texttt{MLKEM768-X25519}. The wrap is secure if \emph{either}
the post-quantum (ML-KEM-768) or the classical (X25519) component holds---hybrid
defense-in-depth against both quantum attack and an implementation flaw in the
young PQ primitive. Recording the sealed DEK conveys nothing without the
corresponding ML-KEM identity.

\textbf{Signatures.} Where the platform authenticates rather than encrypts
(node handshakes, tokens), it uses ML-DSA-65. This closes the asymmetric surface
end-to-end.

\begin{lstlisting}[language=Go,caption={Sealing a per-tenant DEK (HIP-0302).}]
dek := random(32)                 // ChaCha20-Poly1305 key
w, _ := age.Encrypt(out, rcptPQ) // ML-KEM-768 + X25519
w.Write(dek); w.Close()          // sealed DEK -> dek.age
\end{lstlisting}

\section{Surface 1: Native SQLite VFS}
SQLite's virtual file system layer lets a host intercept all I/O. We implement
it in pure Go (\texttt{hanzoai/sqlite3}, a fork of \texttt{ncruces/go-sqlite3};
no cgo) so that a database opened as
\texttt{file:/orgs/acme/app.db?vfs=hanzo} reads and writes encrypted pages
directly, \emph{in process}---no kernel transition. Each page is an AEAD-sealed
block; the database's DEK is the age-wrapped key above. A tenant is therefore one
(or a few per-user) encrypted SQLite files plus their replica; restoring a tenant
is fetching their objects, and revoking one is deleting them.

This makes the relational, key-value, queue, and session needs of an application
collapse into tables inside one per-tenant database---no Postgres, no Redis.

\section{Surface 2: A Dropbox-Class FUSE Mount}
The same substrate is exposed as a POSIX filesystem via a FUSE node
(\texttt{bazil.org/fuse}; a WASI guest mode is planned for sandboxed hosts). The
mount is fully read/write: \texttt{Create}, \texttt{Mkdir}, \texttt{Write},
\texttt{Setattr}, \texttt{Remove}, and \texttt{Fsync} map file operations onto
block \texttt{Put/Get}. Writes are encrypted on the local host \emph{before} a
byte leaves for S3, so the provider sees only ciphertext.

\begin{lstlisting}[language=bash,caption={Mounting S3 as an encrypted folder.}]
vfs mount /mnt/drive \
  --backend s3://bucket/acme?endpoint=... \
  --age-recipient age1pq1... --age-key id.key
\end{lstlisting}

This is the Dropbox-class surface. Three properties make it practical:
\begin{itemize}
\item \textbf{Local-first.} The NVMe cache serves the working set at local
latency; only misses and flushes touch S3. The mount is usable offline for cached
paths and reconciles on reconnect.
\item \textbf{Content-addressed sync.} Because blocks are keyed by content, an
edit re-uploads only changed blocks, and a file's history is a sequence of
immutable block sets---natural versioning and dedup of unchanged regions.
\item \textbf{Sharing by re-wrap.} Sharing a path means re-wrapping its DEK to
another principal's ML-KEM recipient (resolved from \texttt{lux.id}); the bulk
ciphertext is untouched. Revocation rotates the DEK and re-seals forward.
\end{itemize}
Conflicts are handled at the block/inode level: concurrent writers produce
divergent content hashes that are reconciled by a last-writer-wins default with
an explicit conflict copy, mirroring consumer sync semantics while preserving the
encrypted-at-rest invariant.

\section{Replication and Durability}
Durability is HIP-0107: a source adapter streams the substrate's write-ahead
records through the same \texttt{age} envelope to a \texttt{vfs.Backend}, so
SQLite WAL, the file mount, and blockchain state all replicate to S3 through one
pipeline. Restore is hydration of the latest snapshot plus replay of
incrementals; a wiped node rejoins by pulling its tenant objects.

\section{Evaluation}
We benchmark the database surface against a FUSE-mounted equivalent of the same
encrypted substrate (same cipher, same block size), isolating the cost of the
kernel round-trip. All figures are single-tenant on arm64; each tenant is an
independent database, so these are per-tenant, not shared-bottleneck, rates.

\begin{table}[h]\centering\small
\begin{tabular}{lrr}
\toprule
Path & insert (rows/s) & point-read (q/s) \\
\midrule
Native VFS, no crypto & 1{,}250{,}000 & 1{,}130{,}000 \\
Native VFS + AEAD (in-mem) & 980{,}000 & 1{,}100{,}000 \\
\textbf{Native PQ VFS (durable)} & \textbf{720{,}000} & \textbf{59{,}600} \\
FUSE mount (encrypted) & 274{,}000 & 11{,}909 \\
\bottomrule
\end{tabular}
\caption{Native, durable, post-quantum SQLite outperforms the FUSE path by
$2.6\times$ on writes and $5\times$ on point-reads. AEAD encryption costs
$\sim$20\% on writes and $\sim$0 on reads; the dominant FUSE tax is the kernel
transition, not the cryptography.}
\end{table}

The takeaway determines surface selection: \emph{databases} use the native VFS
(no FUSE tax); \emph{files} use the FUSE mount, where POSIX semantics matter more
than raw IOPS and the kernel cost is amortized by the page cache.

\section{Conclusion}
One encrypted, content-addressed, post-quantum block store; two surfaces---an
in-process SQLite VFS for structured tenant state and a read/write FUSE mount for
Dropbox-class files. No second database, no separate sync service, no plaintext
at rest, and no classical-only link in the path from a tenant's bytes to an S3
object. The asymmetric attack surface is ML-KEM and ML-DSA throughout; the
symmetric layer is AEAD. The result is a storage tier that is simultaneously
simpler to operate and durable against a quantum adversary.

\end{document}
